\documentclass[10pt]{article}
\usepackage[a4paper,margin=0.76in]{geometry}
\usepackage[T1]{fontenc}
\usepackage[utf8]{inputenc}
\usepackage{amsmath}
\usepackage{booktabs}
\usepackage{array}
\usepackage{enumitem}
\usepackage{hyperref}
\usepackage{inconsolata}

\newcommand{\SoftCons}{\textsc{SoftCons}}
\newcommand{\HardCons}{\textsc{HardCons}}
\newcommand{\Rev}{\operatorname{Rev}}
\setlength{\parindent}{0pt}
\setlength{\parskip}{4pt}
\setlength{\tabcolsep}{4pt}
\setlist{nosep,leftmargin=*}

\title{\vspace{-1.2cm}DiCo-NLI: Evaluation Plan}
\author{SemEval-2027 Task Proposal Supplementary Material, Item 4}
\date{\vspace{-0.7cm}}

\begin{document}
\maketitle
\vspace{-0.8cm}

\section{Scope}

DiCo-NLI evaluates directional consistency in fine-grained natural language inference.
For each ordered phrase pair, a system predicts exactly one label.
Each track is evaluated independently, and the official result for a system on a track is a three-score profile: weighted $F_1$, \SoftCons{}, and \HardCons{}.
The scores are not collapsed into a single scientific metric because they measure different properties: item-level correctness, directional self-consistency, and paired directional correctness.

The official tracks are English, Spanish, Basque, and Mixed multilingual.
Results will be reported per track.
The task overview may additionally include a monolingual macro-average over English/Spanish/Basque and an all-track macro-average over English/Spanish/Basque/Mixed, but these are secondary summaries and do not replace the per-track scores.

\section{Submission Format and Validation}

For each track, participants submit one prediction file with one row per released evaluation instance.
The required fields are:

\begin{center}
\small
\begin{tabular}{ll}
\toprule
Field & Requirement \\
\midrule
\texttt{instance\_id} & Must match an instance identifier in the evaluation file. \\
\texttt{label} & Must be one of the four official labels. \\
\bottomrule
\end{tabular}
\end{center}

The official label inventory is:
\[
L=\{\texttt{EQUIVALENCE},\texttt{FORWARD\_ENTAILMENT},
\texttt{BACKWARD\_ENTAILMENT},\texttt{NEGATIVE\_OTHER}\}.
\]

The scorer will reject malformed submissions before scoring.
Validation errors include missing instance identifiers, duplicate identifiers, unknown labels, extra predictions for unknown identifiers, empty predictions, invalid encodings, and files whose row count does not match the evaluation set.
The scorer will produce a machine-readable validation report and a human-readable error message.

Final evaluation labels will be hidden during the official evaluation phase.
Gold labels, scorer code, and full result tables will be released after the evaluation phase.

\section{Official Metrics}

\begin{center}
\small
\begin{tabular}{lll}
\toprule
Metric & Unit & Purpose \\
\midrule
Weighted $F_1$ & ordered instances & Item-level label-prediction quality over all four labels. \\
\SoftCons{} & reversible source pairs & Directional compatibility of the two system predictions. \\
\HardCons{} & reversible source pairs & Both directions must be predicted correctly. \\
\bottomrule
\end{tabular}
\end{center}

\subsection{Weighted \texorpdfstring{$F_1$}{F1}}

Weighted $F_1$ is computed over the full label set $L$, including \texttt{NEGATIVE\_OTHER}.
Let $n_\ell$ be the gold support for label $\ell$ and $N=\sum_{\ell\in L}n_\ell$.
For each label $\ell$:
\[
P_\ell=\frac{TP_\ell}{TP_\ell+FP_\ell}, \qquad
R_\ell=\frac{TP_\ell}{TP_\ell+FN_\ell},
\]
\[
F_\ell=
\begin{cases}
\frac{2P_\ell R_\ell}{P_\ell+R_\ell}, & P_\ell+R_\ell>0,\\
0, & P_\ell+R_\ell=0.
\end{cases}
\]
The official item-level score is:
\[
F_w=\sum_{\ell\in L}\frac{n_\ell}{N}F_\ell .
\]
If a precision or recall denominator is zero, the corresponding value is set to zero.
This follows standard zero-division handling and prevents undefined scores.

\subsection{Reversal Operator}

The reversible label set is:
\[
R=\{\texttt{EQUIVALENCE},\texttt{FORWARD\_ENTAILMENT},
\texttt{BACKWARD\_ENTAILMENT}\}.
\]
The deterministic reversal operator is:
\[
\Rev(\texttt{EQUIVALENCE})=\texttt{EQUIVALENCE},
\]
\[
\Rev(\texttt{FORWARD\_ENTAILMENT})=\texttt{BACKWARD\_ENTAILMENT}, \qquad
\Rev(\texttt{BACKWARD\_ENTAILMENT})=\texttt{FORWARD\_ENTAILMENT}.
\]
\texttt{NEGATIVE\_OTHER} is not part of the reversible label algebra.
Gold \texttt{NEGATIVE\_OTHER} instances are included in weighted $F_1$ but excluded from \SoftCons{} and \HardCons{}.

\subsection{\SoftCons{}}

Let $\mathcal{P}$ be the set of reversible source pairs in a track.
For pair $i\in\mathcal{P}$, let $x_i$ be one ordered instance and $x_i^r$ its reversed counterpart.
Let $\hat{y}_i$ and $\hat{y}_i^r$ be the two system predictions.
For each reversible source pair:
\[
s_i =
\begin{cases}
1, & \hat{y}_i\in R,\ \hat{y}_i^r\in R,\ \hat{y}_i^r=\Rev(\hat{y}_i),\\
0, & \text{otherwise.}
\end{cases}
\]
\[
\SoftCons{}=|\mathcal{P}|^{-1}\sum_{i\in\mathcal{P}}s_i .
\]
\SoftCons{} ignores gold correctness and measures whether the two predictions made by the same system are directionally compatible.
For example, predicting \texttt{EQUIVALENCE} in both directions is soft-consistent even if the gold labels are \texttt{FORWARD\_ENTAILMENT} and \texttt{BACKWARD\_ENTAILMENT}.
Predicting \texttt{FORWARD\_ENTAILMENT} in both directions is not soft-consistent.
If either prediction for a reversible pair is \texttt{NEGATIVE\_OTHER}, the pair receives $s_i=0$ because the prediction does not express a reversible directional relation.

\subsection{\HardCons{}}

For the same reversible source pair $i$, let $y_i$ be the gold label for $x_i$.
The gold label for the reversed item is deterministically $\Rev(y_i)$.
For each pair:
\[
h_i =
\begin{cases}
1, & \hat{y}_i=y_i\ \text{and}\ \hat{y}_i^r=\Rev(y_i),\\
0, & \text{otherwise.}
\end{cases}
\]
\[
\HardCons{}=|\mathcal{P}|^{-1}\sum_{i\in\mathcal{P}}h_i .
\]
\HardCons{} is strict paired correctness over complete reversible phrase-pair units.
A pair receives credit only when both directions are correct.
Because the reverse gold label is deterministic, \HardCons{} is always less than or equal to paired item correctness and is expected to be more demanding than ordinary item-level scoring.

\section{Reporting and Leaderboards}

For each submitted run and each track, the scorer will report:

\begin{itemize}
\item the three main scores: weighted $F_1$, \SoftCons{}, and \HardCons{};
\item secondary diagnostics: macro-$F_1$, per-label precision/recall/$F_1$, confusion matrix, number of scored instances, and number of reversible source pairs;
\item validation metadata: accepted/rejected status, missing or duplicate identifiers if any, and timestamped scoring logs.
\end{itemize}

Official tables in the task overview paper will present the three main scores together.
No single scalar will replace them.
If CodaBench requires a display ordering, the sorting convention will be documented in the competition instructions and scorer README as an administrative display choice.
It will not change the definition of the official result, which remains the three-score profile per track.

\section{Evaluation Procedure}

The evaluation pipeline will be frozen before the official evaluation period:
\begin{enumerate}
\item Release trial/training data, submission examples, format checker, scorer documentation, and starter baselines.
\item During development, allow participants to run the scorer on public labeled data and validate submission format locally.
\item During official evaluation, accept submissions on CodaBench against hidden labels.
\item Score accepted submissions with the frozen scorer and hidden gold labels.
\item Publish per-track official scores, secondary diagnostics, and required system metadata.
\item After the evaluation phase, release final gold labels, the scorer, and complete data documentation.
\end{enumerate}

Participants must report use of external data, open-weight or closed/API models, private resources, retrieval systems, prompt engineering, fine-tuning, and ensembling.
This metadata supports task analysis but does not define separate official leaderboards unless announced before the evaluation phase.

\end{document}
